\reportsection{Background: From Retrofit to Security Spine}
\label{sec:phase1}

The first phase of this work compared the existing
\textsf{ark-crypto-primitives} Merkle tree with the Chiesa--Yogev
presentation~\cite{CY26}.  It began as implementation work around open feature
requests: simplify the split hash interface, remove caller-side ordering
footguns, support richer tree shapes, and compress multi-openings.  Read
together, those requests pointed to one broader issue.  The implementation had
many locally reasonable decisions, but no single specification against which
their combined behaviour could be reviewed.

That comparison produced useful engineering results.  In particular,
\code{ImplicitCoPath} performs the book's path-pruning optimisation without
materialising a coordinate bookkeeping layer, and the
\code{Committed}/\code{OpeningProof}/\code{Opening} split makes ordering and
duplicate-index validation part of construction rather than caller discipline.
These results remain relevant, but they are background to the present report.

The stronger finding was methodological.  A reference definition can organise
an API, yet an API alone still does not tell a developer which security claim a
concrete configuration supports.  That is the step from retrofit to security
spine: keep the literature-facing definition, connect it to executable
interfaces, state the security games, derive the parameters, and expose the
evidence boundary.  The technical overview and results show that step for
\textsf{ark-vc} and \textsf{ark-mt}.

\takeaway{Phase 1 supplied the design context and several useful implementation
patterns.  Its lasting role in this report is the lesson that definitions must
remain connected to the code and claims they are meant to govern.}
